\chapter{Arithmetic Foundations}
\label{ch:foundations}

Everything in this thesis rests on three operations: logarithms, percentages, and the normal distribution.
This chapter builds each one from first principles using the Feynman technique.

%% ─────────────────────────────────────────────────────────────────────────────
\section{The Logarithm}

\situation{
You invest \$100. After one year your account shows \$108. After two years, \$116.64.
If someone asks ``how fast did your money grow?'', you could say ``8\%/year'', or you could say ``0.0770/year in log units''. Both answers are correct. The second is more useful for calculations.
}

\problem{
Percentages do not add up correctly across time periods.
A 10\% gain followed by a 10\% loss gives $1.10 \times 0.90 = 0.99$ --- a 1\% net loss, not 0.
Logarithms turn multiplication into addition, making multi-period returns additive.
}

\workedarith{
Day 1 price: 4{,}000. Day 2 price: 4{,}080. Day 3 price: 4{,}050.

\medskip
\textbf{Log return, day 1$\to$2:}
\[
  r_1 = \ln\!\left(\frac{4080}{4000}\right) = \ln(1.02) \approx 0.0198 \approx 1.98\%
\]

\textbf{Log return, day 2$\to$3:}
\[
  r_2 = \ln\!\left(\frac{4050}{4080}\right) \approx -0.0074 \approx -0.74\%
\]

\textbf{Two-day total by addition:} $r_1 + r_2 = 0.0198 - 0.0074 = 0.0124$

\textbf{Direct calculation:} $\ln(4050/4000) \approx 0.0124$ \checkmark

The sum of daily log returns equals the log return over the whole period.
For percentage returns: $(+2\%) + (-0.74\%) = +1.26\%$, but the real answer is $+1.24\%$.
The error is small for small moves, large for large moves.
}

\formaldef{
The natural logarithm $\ln(x)$ is the unique function satisfying:
\[
  e^{\ln(x)} = x, \quad \ln(e^x) = x, \quad \ln(xy) = \ln(x) + \ln(y)
\]
where $e \approx 2.71828$. The log return is defined as:
\[
  r_t = \ln\!\left(\frac{P_t}{P_{t-1}}\right)
\]
The key property: $\sum_{t=1}^T r_t = \ln(P_T/P_0)$ exactly.
}

\analogybreak{
Logs work perfectly for prices, but they require strictly positive inputs.
A price can never be exactly zero (a market would cease trading), so this is rarely binding.
However, log returns of $-100\%$ correspond to $P_t = 0$, which $\ln$ maps to $-\infty$.
In practice, we never observe this.
}

\whyhere{
Every number in this thesis --- every metric, every model input, every evaluation --- is computed from daily log returns $r_t$. We never use raw prices because they are non-stationary (trending over time) and incomparable across markets.
}

%% ─────────────────────────────────────────────────────────────────────────────
\section{Mean and Standard Deviation}

\situation{
You measure the temperature every day for a week: 15, 17, 14, 19, 16, 21, 13 degrees.
You want two summaries: the ``typical'' value, and how much the values vary.
}

\problem{
One number never tells the whole story. Knowing the average temperature of 16.4\textdegree C does not tell you that some days hit 21\textdegree. We need both location and spread.
}

\workedarith{
Seven values: $x = \{15, 17, 14, 19, 16, 21, 13\}$.

\medskip
\textbf{Mean:} $\bar{x} = (15+17+14+19+16+21+13)/7 = 115/7 \approx 16.43$

\textbf{Deviations:} $\{-1.43, +0.57, -2.43, +2.57, -0.43, +4.57, -3.43\}$

\textbf{Squared deviations:} $\{2.04, 0.33, 5.90, 6.60, 0.18, 20.87, 11.77\}$

\textbf{Variance:} $s^2 = (2.04+0.33+\ldots)/6 = 47.69/6 = 7.95$ (divide by $n-1$ for sample)

\textbf{Standard deviation:} $s = \sqrt{7.95} \approx 2.82$\textdegree C
}

\formaldef{
For a sample $\{x_1, \ldots, x_n\}$:
\[
  \bar{x} = \frac{1}{n}\sum_{i=1}^n x_i, \qquad
  s^2 = \frac{1}{n-1}\sum_{i=1}^n (x_i - \bar{x})^2, \qquad
  s = \sqrt{s^2}
\]
The $(n-1)$ denominator corrects for the fact that we estimated $\bar{x}$ from the same sample (Bessel's correction). The population versions use $n$.
}

\analogybreak{
Standard deviation describes spread around the mean, but only makes full sense when the distribution is symmetric. For skewed distributions (like stock returns, which can fall much further than they rise), standard deviation understates downside risk.
}

\whyhere{
Daily volatility $\sigma = \text{sd}(r_t)$ is the foundational measure in all five markets. In our 20-year data: MOEX $\sigma=1.87\%$/day, BOVESPA $1.63\%$. Annualised: multiply by $\sqrt{252}$. Standard deviation also enters GARCH, VaR, discriminative AUC, and every distance metric we compute.
}

%% ─────────────────────────────────────────────────────────────────────────────
\section{The Normal Distribution and Why It Fails for Markets}

\situation{
Roll a fair die 100 times and record the average result. Do this experiment 10{,}000 times. Plot the 10{,}000 averages. They form a perfect bell shape: most near 3.5, rarely below 2.5 or above 4.5. This is the Central Limit Theorem in action.
}

\problem{
The normal distribution makes precise predictions about extreme events. For a bell curve, a five-standard-deviation event should occur once every 6{,}922 years. We observed 76 such events across 24{,}758 BRICS market-days. The model is not slightly wrong; it is orders-of-magnitude wrong.
}

\workedarith{
The Gaussian probability of $|x| > 5\sigma$ is:
\[
  \Prob(|Z| > 5) = 2 \times \Phi(-5) \approx 5.73 \times 10^{-7}
\]
With 252 trading days per year, the expected waiting time is:
\[
  \text{Wait} = \frac{1}{252 \times 5.73 \times 10^{-7}} \approx 6{,}922 \text{ years}
\]
We have 24{,}758 market-days across 5 markets $\approx 100$ market-years total.
The Gaussian predicts $100 \times 252 / 6{,}922 \approx 0.014$ such days.
We observed 76. The ratio is $76/0.014 \approx 5{,}400$.
}

\formaldef{
The normal distribution $\N(\mu, \sigma^2)$ has PDF:
\[
  f(x) = \frac{1}{\sigma\sqrt{2\pi}}\exp\!\left(-\frac{(x-\mu)^2}{2\sigma^2}\right)
\]
The \textbf{excess kurtosis} measures tail weight relative to the normal:
\[
  \text{Kurt} = \frac{\E[(X-\mu)^4]}{(\E[(X-\mu)^2])^2} - 3
\]
Normal: $\text{Kurt}=0$. Heavy-tailed distributions: $\text{Kurt}>0$. BRICS markets: 5.73--65.70.

\medskip
\textbf{Important caveat:} kurtosis requires $\E[X^4] < \infty$. For power-law tails $P(|X|>x)\sim x^{-\alpha}$, the $k$-th moment exists only if $k < \alpha$. With BRICS tail indices $\hat{\alpha}=2.54$--$3.21$, the fourth moment is infinite everywhere and the third moment is infinite in 4 of 5 markets. Sample kurtosis diverges with $n$ rather than converging.
}

\analogybreak{
The Central Limit Theorem says averages of independent identically distributed random variables converge to a normal. Market returns are \emph{not} independent (volatility clusters), and their distribution is \emph{not} finite-variance when $\alpha < 2$ (our markets have $\alpha>2$, so variance is finite, but fourth moments are not). CLT still applies to averages, but individual extreme events are not well-described by the normal.
}

\whyhere{
The failure of the normal distribution is the entire motivation for this thesis. If market returns were Gaussian, banks would be adequately protected by Gaussian VaR models. Because they are not, we need generative models that reproduce heavy tails and volatility clustering --- and a rigorous evaluation framework to verify they do.
}

%% ─────────────────────────────────────────────────────────────────────────────
\section{Conditional Probability}

\situation{
A city has 60\% chance of rain on any day. But if it rained yesterday, the chance of rain today is 80\%. Knowing yesterday's outcome changes the probability of today's.
}

\workedarith{
Let $A$ = ``big market move today'', $B$ = ``big market move yesterday''.

MOEX 20-year data: $\Prob(A) = 3.8\%$, $\Prob(A|B) = 22.9\%$.

Ratio: $\Prob(A|B)/\Prob(A) = 22.9/3.8 = 6.1$.

This means: knowing that yesterday was a big-move day makes today's big-move probability $\mathbf{6.1\times}$ higher than the unconditional baseline.
}

\formaldef{
\[
  \Prob(A \mid B) = \frac{\Prob(A \cap B)}{\Prob(B)}
\]
Events $A$ and $B$ are \textbf{independent} if $\Prob(A|B) = \Prob(A)$.
Market returns are not independent: $\Prob(A|B) \neq \Prob(A)$.
The ratio $\Prob(A|B)/\Prob(A) = 6.1$ measures the strength of \textbf{volatility clustering}.
}

\analogybreak{
Conditional probability does not imply causality. MOEX big moves after invasion news are driven by the news, not by the prior day's move per se. Correlation in extremes is not the same as predictability of direction.
}

\whyhere{
The shuffled-control experiment destroys conditional structure: $\Prob(A|B)_\text{shuffled} \approx \Prob(A)_\text{shuffled}$ after permutation. Any generator that reproduces only the marginal distribution $\Prob(A)$ but not the conditional $\Prob(A|B)$ will pass distributional tests but fail temporal ones.
}
